\section{Direct embedding lemmas}
\label{sec:direct}

Neighborhood transfer will handle edges between a small
high-degree set and its complement. Two direct estimates cover the other parts:
one for graphs with bounded maximum degree and one for small
star forests. Neither requires membership in the typical family.

\subsection{A maximum-degree bound}

\begin{proposition}\label{prop:degree-embedding}
There is an absolute integer $N_2$ such that, for $n\geq N_2$
and $r,s$ satisfying \eqref{eq:parameters}, every graph $J$
with
\[
 |V_+(J)|\leq n/4,\qquad \degen(J)\leq s,
 \qquad\Delta(J)\leq n^{2/3}
\]
and at most $n$ vertices satisfies $p_c(J;n)\leq r$.
\end{proposition}
\begin{proof}
Discard isolated vertices and assume that an edge remains.
We partition the vertices into layers that are independent sets, constructed
in reverse order. In any residual subgraph with $m'$ vertices,
the average degree is at most $2s$. At least $m'/2$ vertices
therefore have degree at most $4s$. The graph they induce has
an independent set of size at least $m'/(2(4s+1))$, by greedily
choosing a vertex and deleting it and its neighbors. Take such
a set as the last remaining layer and repeat on the residual
graph. Geometric decrease gives
\begin{equation}\label{eq:layers}
 t\leq C s\log n\leq C'(\log n)^2
\end{equation}
layers. Every vertex has at most $4s$ neighbors in earlier layers.

For each layer choose a disjoint bin in $[n]$ of size equal to
the layer size plus
\[
 b=\left\lfloor\frac{n}{4t}\right\rfloor.
\]
The bins fit, since their total size is at most $n/4+tb\leq n/2$.
Embed the layers in order. At a given step expose the edges
between the current bin and the earlier images, and seek a
matching of the layer into its bin. These edges are fresh
independent Bernoulli-$r$ variables conditional on the entire
past: all previous exposures were between earlier bins.

Write $\Delta=\Delta(J)$. Fix a subset $U$ of the current
layer with $|U|=u\geq1$. Each vertex of $U$ requires adjacency
to at most $4s$ earlier images. Each earlier image occurs in at
most $\Delta$ requirements, so the conflict graph joining
intersecting requirements has maximum degree at most $4s\Delta$.
It has an independent set of size at least
$u/(1+4s\Delta)$, whose requirements are pairwise disjoint.

For a set $W$ of $u-1$ vertices in the current bin, at least $b$
vertices of the bin remain outside $W$. The probability that
none is usable for any vertex of $U$ is at most
\begin{equation}\label{eq:hall-failure}
 \exp\!\left(-\frac{bu r^{4s}}{1+4s\Delta}\right).
\end{equation}
Indeed, restrict to the disjoint requirements just selected.
Each succeeds at a given candidate with probability at least
$r^{4s}$, independently over those requirements and over
candidates. If an empty requirement is present, the failure
event is impossible, so the bound remains valid.

Uniformly under \eqref{eq:parameters} and \eqref{eq:layers},
for sufficiently large $n$,
\begin{equation}\label{eq:hall-exponent}
 \frac{b r^{4s}}{1+4s\Delta}
 \geq c\frac{n^{1-4/100-2/3}}{(\log n)^3}
 =c\frac{n^{22/75}}{(\log n)^3}\geq5\log n.
\end{equation}
Here $b\geq c'n/(\log n)^2$ and $r^{4s}\geq n^{-4/100}$.
There are at most $n^{2u}$ choices of $U,W$ for each $u$.
Failure of Hall's condition implies such a pair with all usable
images of $U$ in $W$. By \eqref{eq:hall-failure}--\eqref{eq:hall-exponent},
the conditional failure probability at a layer is at most
\[
 \sum_{u\geq1}n^{2u}e^{-5u\log n}
 =\sum_{u\geq1}n^{-3u}=O(n^{-3}).
\]
Summing this bound over the $O((\log n)^2)$ layers shows that
the embedding succeeds with probability tending to one,
uniformly over the allowed parameters. Taking an absolute
$N_2$ so that it succeeds with probability at least $1/2$
proves the proposition.
\end{proof}

\subsection{Small star forests}

A star forest is a graph whose nontrivial components are stars.
The next argument adapts the fresh-center embedding in
\cite[Claim~6.6]{AHPT26}. Small support is imposed directly,
so it is unnecessary to divide the stars by their degrees in
advance.

\begin{proposition}\label{prop:stars}
Put $A=400e^4$. There is an absolute integer $N_3$ such that,
for every $n\geq N_3$ and every star forest $J$ on at most $n$
vertices with $|V_+(J)|\leq n/100$,
\begin{equation}\label{eq:star-threshold}
 p_c(J;n)\leq2Aq(J;n).
\end{equation}
If $q(J;n)<1/(2A)$, the stronger bound
$p_c(J;n)\leq Aq(J;n)$ holds.
\end{proposition}
\begin{proof}
The edgeless case is immediate. Put $q_J=q(J;n)>0$.
If $q_J\geq1/(2A)$, then $2Aq_J\geq1$, so
\eqref{eq:star-threshold} is automatic. We may therefore
assume $q_J<1/(2A)$ and set $p=Aq_J\in(0,1/2)$.

For $\ell\geq1$, let $m_\ell$ be the number of components
of $J$ with exactly $\ell$ leaves, ignoring isolated vertices.
If $m=m_\ell>0$, the copies on $[n]$ of $m$ disjoint
$\ell$-edge stars cover $\calF_J$. Their number is at most
\[
 C_\ell=\frac{n^{m(\ell+1)}}{(\ell!)^m m!}.
\]
This follows by counting vertex injections and dividing by
permutations of the leaves and of the stars. For $\ell=1$
there are additional automorphisms, which only decrease the
actual count. Since such a star forest exists on at most $n$
vertices, $C_\ell\geq1$. Using its cover at
$(1/(2C_\ell))^{1/(m\ell)}\in(0,1)$ yields
\begin{equation}\label{eq:star-count}
 q_J^\ell\geq\frac{\ell!}{n^{\ell+1}}(m!/2)^{1/m}
 \geq\frac{\ell!}{n^{\ell+1}}\frac{m_\ell}{2e}.
\end{equation}
The last step uses $m!\geq(m/e)^m$ and $2^{1/m}\leq2$.

Partition $[n]$ into a center bin of size $\lfloor n/2\rfloor$
and a leaf bin containing the remaining vertices. Process the
components of $J$ one at a time. For a star with $\ell$ leaves,
test unused centers successively until one has at least $\ell$
neighbors among the unused leaf vertices. Use that center and
any $\ell$ such neighbors, and never test that center again.
Failed centers are also discarded.

Throughout the process, at least $b'=\lfloor n/3\rfloor$
leaf vertices remain. The unused leaf set is determined by
previously tested center rows, and a fresh center row has not
been exposed. Thus each trial, conditional on the past, has
success probability at least
\begin{equation}\label{eq:star-success}
 u_\ell=\Prob(\Bin(b',p)\geq\ell).
\end{equation}
For the purpose of counting trials, append infinitely many
auxiliary centers with fresh independent rows to the same leaf
bin. This enlarged process agrees with the actual process until
the original center bin is exhausted. The number $T$ of trials
needed to finish all stars in the enlarged process satisfies
\begin{equation}\label{eq:expected-trials}
 \E T\leq\sum_{\ell:m_\ell>0}\frac{m_\ell}{u_\ell},
\end{equation}
because the conditional expected waiting time for each success
is at most $1/u_\ell$.

If $\mu=b'p\geq2\ell$, the second-moment bound
$\E[\Bin(b',p)^2]\leq\mu^2+\mu$ and the Paley--Zygmund
inequality give
\[
 u_\ell\geq\frac14\frac{\mu^2}{\mu^2+\mu}\geq\frac16,
\]
since $\mu\geq2$. Otherwise $b'p<2\ell$. For every occupied
star size, $\ell\leq n/100$, and for large $n$,
$b'-\ell+1\geq n/4$. As $p<1/2$, the inequality
$\log(1-p)\geq-2p$ implies
\begin{align}
 u_\ell
 &\geq\binom{b'}{\ell}p^\ell(1-p)^{b'-\ell}
 \geq\frac{(n/4)^\ell p^\ell}{\ell!}e^{-4\ell}\notag\\
 &\geq\left(\frac{A}{4e^4}\right)^\ell
        \frac{m_\ell}{2e n}
 =100^\ell\frac{m_\ell}{2e n},
 \label{eq:star-small-mean}
\end{align}
where the last inequality uses \eqref{eq:star-count}.

There are at most $n/200$ nontrivial components. Combining
the two bounds in \eqref{eq:expected-trials} gives
\[
 \E T\leq6\frac{n}{200}+2e n\sum_{\ell\geq1}100^{-\ell}
 =\left(\frac3{100}+\frac{2e}{99}\right)n<\frac n8.
\]
For large $n$, Markov's inequality shows that
$T\leq\lfloor n/2\rfloor$ with probability at least $1/2$.
On this event all centers used belong to the actual random graph,
so the construction embeds $J$ in $G(n,p)$. This proves
$p_c(J;n)\leq Aq_J$ in the remaining case.
\end{proof}

\section{Constructing the fixed partition}
\label{sec:assembly}

We now prove Theorem~\ref{thm:main}. The high-degeneracy case
first uses the typical-family cover and isolates a small set $Y$
so that every vertex outside $Y$ has few neighbors in $Y$ and
bounded total degree. Neighborhood transfer handles edges across
this cut, while the maximum-degree estimate handles edges inside
and outside $Y$ after further splitting. The bounded-degeneracy
case uses only small star forests. The partition choices in both
cases are deterministic.

\begin{proof}[Proof of Theorem~\ref{thm:main}]
Fix
\begin{equation}\label{eq:constants}
 M=B=1000,\qquad k=2K_0M(B+1)^2.
\end{equation}
Choose an absolute integer $n_0\geq2$ large enough for
Propositions~\ref{prop:transfer}, \ref{prop:degree-embedding},
and~\ref{prop:stars}, and so that, for every $n\geq n_0$,
\begin{equation}\label{eq:deterministic-cutoffs}
 \frac{8\log n}{\log2}\sqrt n\leq n^{2/3},\qquad
 n^{2/3}\leq n/4,\qquad
 2\lceil n/B\rceil\leq n/100.
\end{equation}
All these conditions have absolute cutoffs. Set
\[
 L=\max\{2A,64,n_0^2\}.
\]
If $H$ is edgeless, use only empty pieces. Otherwise write
$q=q(H;n)>0$ and $d=\degen(H)\geq1$.
When $n<n_0$, Lemma~\ref{lem:q-lower} gives
$Lq\geq n_0^2/(n(n-1))\geq1$. When $q\geq1/(2A)$,
we also have $Lq\geq1$. In either case put all edges in one
piece and use empty pieces for the rest. Henceforth
\begin{equation}\label{eq:nontrivial-range}
 n\geq n_0,\qquad q<1/(2A).
\end{equation}

\paragraph{Large degeneracy: $d\geq M$.}
Set
\[
 r=2q,\qquad s=\lceil2d/M\rceil.
\]
By Lemma~\ref{lem:q-lower}, $r\geq n^{-2/d}$, whereas
\eqref{eq:nontrivial-range} gives $r<1/A<1/2$. Consequently
\begin{equation}\label{eq:d-upper}
 d<\frac{2\log n}{\log2},\qquad s\leq3d/M,
 \qquad r^s\geq n^{-6/M}\geq n^{-1/100}.
\end{equation}
These bounds also give $s\leq10\log n$, so
\eqref{eq:parameters} holds.

Take the family $\calD$ from Proposition~\ref{prop:transfer}.
Corollary~\ref{cor:typical-cover} supplies a copy of $H$ covered
by $K_0$ members $G_1,\ldots,G_{K_0}$ of $\calD$.
Fix those graphs and the labelled copy, and use this copy to
label $H$ on a subset of $[n]$. Assign each edge of $H$ to
one $G_i$ containing it. This gives $K_0$ initial edge pieces,
each contained in its specified typical graph.

We separate the high-degree vertices, following the obstacle
identified in \cite[Section~8]{AHPT26}. Initially let
\[
 Y_0=\{v\in V(H):\deg_H(v)>\sqrt n\}.
\]
By \eqref{eq:edge-bound}, $|Y_0|\leq2d\sqrt n$.
Starting from $Y_0$, add an outside vertex whenever it has more
than $2d$ neighbors in the current set, and stop when no such
vertex remains. Let $Y$ be the resulting set and let
$a=|Y\setminus Y_0|$. Each added vertex contributes more than
$2d$ edges to earlier vertices of the set; no edge is counted
twice. Thus
\[
 2da\leq|E(H[Y])|\leq d(|Y_0|+a),
\]
giving $|Y|\leq2|Y_0|\leq4d\sqrt n\leq n^{2/3}$ by
\eqref{eq:d-upper} and \eqref{eq:deterministic-cutoffs}.
The conclusion also holds if $Y_0$ is empty, in which case no
vertex is added. At termination every $w\notin Y$ satisfies
\begin{equation}\label{eq:outside-bounds}
 \deg_H(w)\leq\sqrt n,\qquad |N_H(w)\cap Y|\leq2d.
\end{equation}

Partition $V(H)\setminus Y$ into $B$ possibly empty blocks
of size at most $\lceil n/B\rceil$. Subdivide each initial
edge piece as follows.

\emph{Edges inside $Y$.}
Apply Lemma~\ref{lem:split} with $m=M$ to obtain $M$ pieces
of degeneracy at most $\lceil d/M\rceil\leq s$.
Each is supported on at most $|Y|\leq n^{2/3}\leq n/4$
vertices and has maximum degree at most $n^{2/3}$.
Proposition~\ref{prop:degree-embedding} gives threshold at most $r$.

\emph{Edges outside $Y$.}
First partition by the unordered pair of endpoint blocks,
allowing both endpoints in the same block. There are
$B(B+1)/2$ such pairs. Apply Lemma~\ref{lem:split} with
$m=M$ to each resulting graph. The pieces have degeneracy
at most $s$, maximum degree at most $\sqrt n$, and support
size at most $2\lceil n/B\rceil\leq n/100$.
Proposition~\ref{prop:degree-embedding} again gives threshold
at most $r$.

\emph{Edges between $Y$ and its complement.}
First partition by the block of the endpoint outside $Y$.
At each such endpoint, distribute its edges in the initial
piece among $M$ groups of size at most $\lceil2d/M\rceil=s$,
using \eqref{eq:outside-bounds}. This assigns every cut edge
once. Each resulting bipartite piece is still contained in its
specified $G_i\in\calD$, has at most $\lceil n/B\rceil\leq n/8$
positive-degree vertices outside $Y$, and has all those degrees
at most $s$. Proposition~\ref{prop:transfer} gives threshold
at most $r_*\leq32r=64q$.

All pieces therefore have threshold at most $64q\leq Lq$.
Their number is at most
\[
 K_0M\left(1+\frac{B(B+1)}2+B\right)
 =\frac{K_0M(B+1)(B+2)}2\leq k.
\]

\paragraph{Bounded degeneracy: $1\leq d<M$.}
Lemma~\ref{lem:split}, with $m=d$, partitions $E(H)$ into
$d$ graphs of degeneracy at most one, hence into forests.
Indeed, the last vertex of a cycle in the certifying order
would have two preceding neighbors.

Root every nontrivial tree. Color each parent--child edge by
the parity of the depth of its parent. Each color class is a
star forest: its centers occur at one depth parity, and their
children have the other parity and no other incident edge of
that color. This is the parity splitting used in
\cite[Claim~6.2]{AHPT26}.

Partition $V(H)$ into $B$ blocks of size at most
$\lceil n/B\rceil$, and split each star forest by unordered
endpoint-block pairs. Every resulting piece $J$ remains a star
forest and has
$|V_+(J)|\leq2\lceil n/B\rceil\leq n/100$.
By Proposition~\ref{prop:stars} and \eqref{eq:monotonicity},
\[
 p_c(J;n)\leq2Aq(J;n)\leq2Aq(H;n)\leq Lq.
\]
The number of pieces is at most
$2d\,B(B+1)/2=dB(B+1)\leq k$.

Finally, add empty pieces if necessary. In the large-degeneracy
case, pull the labelled edge pieces back to the original $H$.
The typical graphs, the copy of $H$, all vertex orders, and
all assignments have been fixed using only $H,n$; none depends
on the random host used to test a piece's threshold. This
completes the proof.
\end{proof}

\bibliographystyle{amsplain}
\bibliography{references}
\end{document}
